% Irodov Part IV §4.1 independent original problems 21--33.
% User source PDF pp.165--166; printed pp.156--157.
\oscproblem{21}{156}
单摆钟安置于电梯中。电梯从静止开始以恒加速度
$w$ 向上运动（$w<g$），升高高度 $h$ 时将加速度方向反转而大小仍为 $w$。
问从开始运动起经过多久，钟才重新显示正确时间？
\exsolution
上升加速阶段等效重力
$g_+=g+w$，
摆钟相对于真实时间的走时速率为
$r_+=\sqrt{(g+w)/g}>1$。
从静止上升高度 $h$ 所用真实时间
$t_1=\sqrt{2h/w}$。
反向加速阶段等效重力
$g_-=g-w$，
走时速率
$r_-=\sqrt{(g-w)/g}<1$。
要让累计超前量被追平，
$(r_+-1)t_1+(r_--1)t_2=0$，
所以
$t_2=(r_+-1)t_1/(1-r_-)$。
令 $\eta=w/g$，
\[
\boxed{t=t_1+t_2=
\sqrt{\frac{2h}{w}}\,
\frac{\sqrt{1+\eta}-\sqrt{1-\eta}}
{1-\sqrt{1-\eta}}}.
\]
\excheck 该时刻钟再次准点发生在加速度反向之后；所求为真实经过时间，不是表盘累计时间。

\oscproblem{22}{156}
图4.2密度计（浮计）质量
$m=50$g，细管半径
$r=3.2$mm，浸在液体密度
$\rho=1.00\,\mathrm{g/cm^3}$ 中。
将它沿竖直方向轻微压下后释放，忽略液体阻力，求小振动周期。
\exsolution
下移 $x$ 后额外排开液体体积为
$\Delta V=\pi r^2x$，
增量浮力
$\Delta F_b=\rho g\pi r^2x$，
相对于静平衡的回复力为
$F_x=-\rho g\pi r^2x$。
等效弹性系数
$k=\rho g\pi r^2$，
故
\[
\boxed{T=2\pi\sqrt{\frac{m}{\rho g\pi r^2}}}
=\boxed{2.5\,\mathrm s}\quad\text{（约）}.
\]
\excheck 必须使用浮计\emph{细管}的截面积，
不是浮体整体横截面积。

\oscproblem{23}{156}
一根未变形弹簧两端均固定，
其整体弹性系数为
$k=3.2\,\mathrm{N/m}$。
在距一端为全长
$\eta=1/3$ 的位置附加质量
$m=25$g 的小物体。
忽略重力和弹簧质量，求小物体沿弹簧的微小纵向振动周期。
\exsolution
对同一均匀弹簧材料，
弹簧刚度与其长度成反比。
分段后左段刚度
$k_1=k/\eta$，
右段
$k_2=k/(1-\eta)$。
质量向右移动 $x$ 时，
左段被拉长、右段被压缩，两段均提供回复力，
所以
\[
k_{\rm eff}=k_1+k_2
=\frac{k}{\eta(1-\eta)}.
\]
因此
\[
\boxed{T=2\pi\sqrt{\frac{m\eta(1-\eta)}{k}}
\simeq0.26\,\mathrm s}.
\]
\excheck 本版答案页给出\emph{同一公式}，
但数值标成 $0.13$s；代
$m=0.025$kg、$\eta=1/3$、$k=3.2$N/m
得 $0.26$s，数值相差约2倍，已列勘误。

\oscproblem{24}{156}
图4.3：质量为 $m$ 的小物体在水平光滑轨道上，
左右分别通过刚度为 $k_1,k_2$ 的轻弹簧连接固定支座。
求微小纵向振动周期。
\exsolution
相对平衡位置向右位移 $x$ 时，
左弹簧伸长 $x$、右弹簧缩短 $x$，
两者回复力同向，合力
$F=-(k_1+k_2)x$。
运动方程
$m\ddot x+(k_1+k_2)x=0$，
故
\[
\boxed{T=2\pi\sqrt{\frac{m}{k_1+k_2}}}.
\]
\excheck 两端各连接固定墙的两只弹簧对移动质量而言属于刚度相加，不是串联等效。

\oscproblem{25}{156}
图4.4：一质量为 $m$ 的物体悬挂在\emph{首尾串联}的两根轻弹簧下，
两弹簧刚度分别为 $k_1,k_2$。
求竖直小振动周期。
\exsolution
两弹簧串联时力相同，
总位移 $x=x_1+x_2=F/k_1+F/k_2$，
等效刚度
$k_{\rm eq}=(1/k_1+1/k_2)^{-1}
=k_1k_2/(k_1+k_2)$。
竖直重力只移动静平衡位置，不改变小振动曲率，故
\[
\boxed{T=2\pi\sqrt{\frac{m(k_1+k_2)}{k_1k_2}}}.
\]
\excheck 与上一题不同，图4.4中弹簧确实首尾串联，不能把刚度直接相加。

\oscproblem{26}{156}
质量为 $m$ 的小物体固定在长 $2\ell$ 的水平张紧细绳中点，
绳的静张力为 $T_0$。
忽略重力和绳质量，求横向小振动角频率。
\exsolution
质点离开中点横向位移 $y$，
两侧绳段各长约 $\ell$，
小角度下
$\sin\alpha\simeq y/\ell$，
两侧绳张力的横向合力
$F_y=-2T_0y/\ell$。
所以
$m\ddot y+(2T_0/\ell)y=0$，
\[
\boxed{\omega_0=\sqrt{\frac{2T_0}{m\ell}}}.
\]
\excheck 两段绳各提供一个横向分力，所以刚度带系数2。

\oscproblem{27}{156}
图4.5弯管内装水银总质量
$m=200$g，等截面积
$S=0.50\,\mathrm{cm^2}$，
右臂与竖直夹角
$\theta=30^\circ$。
忽略黏滞，求水银振荡周期。
\exsolution
沿管路令水银柱整体偏移距离 $x$：
左臂液面升高 $x$，
右臂液面沿斜管下降 $x$，
竖直高度改变为 $x\cos\theta$。
两端自由液面竖直高度差为
$(1+\cos\theta)x$，
由液压产生的合回复力
$F=-\rho gS(1+\cos\theta)x$。
整个水银柱质量为 $m$，故
\[
\boxed{T=2\pi
\sqrt{\frac{m}{\rho gS(1+\cos\theta)}}}.
\]
取水银密度 $\rho\simeq13.6\times10^3\,\mathrm{kg/m^3}$，
得
$\boxed{T\simeq0.80\,\mathrm s}$。
\excheck $S$ 的单位换算是
$0.50\,\mathrm{cm^2}=5.0\times10^{-5}\,\mathrm{m^2}$。

\oscproblem{28}{156--157}
图4.6中，一根均匀杆架在两只以相反方向旋转的轮子上。
两轮轴间距 $\ell=20$cm，杆—轮间滑动摩擦系数为
$k=0.18$。
证明杆作简谐振动并求周期。
\exsolution
令杆质心相对两轮中点的位移为 $x$，
左右法向支持力满足
$N_1+N_2=mg$，
力矩平衡
$N_2-N_1=2mgx/\ell$，
所以
$N_1=mg(1/2-x/\ell)$、
$N_2=mg(1/2+x/\ell)$。
图中两轮顶面分别向中心移动，
摩擦力分别为右向 $kN_1$ 与左向 $kN_2$，
合力
$F=k(N_1-N_2)=-2kmgx/\ell$。
因此
\[
\ddot x+\frac{2kg}{\ell}x=0,
\qquad
\boxed{T=2\pi\sqrt{\frac{\ell}{2kg}}
=\pi\sqrt{\frac{2\ell}{kg}}
\simeq1.5\,\mathrm s}.
\]
\excheck 需保持杆始终压在两轮上，
即 $|x|<\ell/2$，此时两支持力均为正。

\oscproblem{29}{157}
设地球为密度均匀的球体，
从北极到南极沿地轴贯穿一条直隧道，
忽略空气阻力和地球自转影响。
求 (a) 质点在隧道中的运动方程；
(b) 从一极到另一极的时间；
(c) 通过地心的速率。
\exsolution
球对称均匀质量分布使距离地心 $x$ 处的引力加速度
$g(x)=-gx/R$，
其中 $R$ 为地球半径、
$g$ 为地表重力加速度。
因此
\[
\boxed{\ddot x+\frac gR x=0}.
\]
频率 $\omega_0=\sqrt{g/R}$，
从一端到另一端是半周期：
\[
\boxed{t_{\rm pole}=\pi\sqrt{\frac Rg}
\simeq42\,\mathrm{min}}.
\]
从地表 $x=R$ 静止释放，地心最大速度
\[
\boxed{v_{\max}=\omega_0R=\sqrt{gR}
\simeq7.9\,\mathrm{km/s}}.
\]
\excheck 原题是沿地球\emph{自转轴}的隧道，
所以忽略自转的离心效应与此几何一致。

\oscproblem{30}{157}
图4.7：长 $\ell$ 的单摆悬点 O 相对地面以恒加速度矢量
$\bm w$ 作任意方向运动，
与重力加速度矢量 $\bm g$ 夹角为 $\beta$。
求小振动周期。计算
$\ell=21$cm、$w=g/2$、$\beta=120^\circ$ 时的周期。
\exsolution
在随悬点加速的非惯性参考系，
有效重力为
$\bm g_{\rm eff}=\bm g-\bm w$，
模
\[
g_{\rm eff}=\sqrt{g^2+w^2-2gw\cos\beta}.
\]
摆球围绕有效重力方向静平衡，小角运动
$\ddot\theta+(g_{\rm eff}/\ell)\theta=0$，
故
\[
\boxed{T=2\pi\sqrt{\frac{\ell}
{\sqrt{g^2+w^2-2gw\cos\beta}}}}.
\]
代 $w=g/2$、$\cos120^\circ=-1/2$，
$g_{\rm eff}=g\sqrt{7/4}$，
得
$\boxed{T\simeq0.80\,\mathrm s}$。
\excheck 不能直接把加速度的模 $w$ 与 $g$ 相减；
需作矢量差。

\oscproblem{31}{157}
图4.8：质量 $m=0.20$kg 的滑套 M 位于两只相同弹簧之间，
两弹簧总有效刚度为 $k=20$N/m，
可沿水平杆 AB 无摩擦滑动。
整个装置绕通过杆中点的竖直轴以恒角速度
$\Omega=4.4\,\mathrm{rad/s}$ 旋转。
求小振动周期，以及使振动不可能出现的 $\Omega$ 范围。
\exsolution
在旋转参考系，
离轴位移 $x$ 的离心力为 $m\Omega^2x$，
两弹簧回复力为 $-kx$，故
\[
m\ddot x=-(k-m\Omega^2)x,
\quad
\omega_{\rm eff}^2=\frac km-\Omega^2.
\]
只在 $k/m>\Omega^2$ 时稳定：
\[
\boxed{T=\frac{2\pi}{\sqrt{k/m-\Omega^2}}
\simeq0.70\,\mathrm s}.
\]
当
$\boxed{|\Omega|\ge\sqrt{k/m}=10\,\mathrm{rad/s}}$
时平衡点不再具有正回复曲率，
不存在该中心平衡附近的小谐振动。
\excheck 阈值等号时频率为零，是临界状态而不是有限周期振动。

\oscproblem{32}{157}
木块置于水平振动木板上，木板作简谐运动，
振幅 $a=10$cm、周期 $T=1.0$s。
当木块刚开始相对木板滑动时，求静摩擦系数的临界值。
\exsolution
木板最大水平加速度为
$a\omega^2$，
$\omega=2\pi/T$。
木块不滑动所需最大水平力为
$ma\omega^2$，
最大静摩擦力为 $\mu mg$。
临界条件
$\mu mg=ma\omega^2$，所以
\[
\boxed{\mu_{\rm crit}=\frac{4\pi^2a}{gT^2}
\simeq0.40}.
\]
\excheck 当 $\mu<\mu_{\rm crit}$ 时开始滑动；
原题英文对幅度与周期的语序有些混乱，
本稿以给出的 $a,T$ 与书末公式为准。

\oscproblem{33}{157}
长 $\ell=80$cm 的单摆在小角条件下，
求三种初始条件对应角位移：
(a) 初偏角 $3.0^\circ$，无初速；
(b) 初在平衡处，摆球获得水平初速 $0.22$m/s；
(c) 原版题面印作初偏角 $30^\circ$，
并获得大小 $0.22$m/s、指向平衡位置的速度。
\exsolution
小角单摆角频率
$\omega_0=\sqrt{g/\ell}=3.5$rad/s，
一般解为
$\theta(t)=\theta_0\cos\omega_0t+
[\dot\theta_0/\omega_0]\sin\omega_0t$。
(a)
\[
\boxed{\theta(t)=3.0^\circ\cos(3.5t)}.
\]
(b) $\dot\theta_0=v_0/\ell=0.275$rad/s，
$|\dot\theta_0|/\omega_0=0.0786$rad$=4.5^\circ$，
所以
\[
\boxed{\theta(t)=4.5^\circ\sin(3.5t)}.
\]
(c) 以印刷题面\emph{字面}初偏角
$30^\circ$ 代入线性化公式得到
\[
\boxed{\theta(t)\approx30^\circ\cos(3.5t)
-4.5^\circ\sin(3.5t)}.
\]
但 $30^\circ$ 已非严格小角，若求精确摆动应解
$\ddot\theta+(g/\ell)\sin\theta=0$。
书末答案却给
$5.4^\circ\cos(3.5t+1.0)$，
它对应初偏角约 $3.0^\circ$
（而非 $30^\circ$），属于明显的原版题面/答案歧义。
\excheck 不把 $30^\circ$ 擅改为 $3.0^\circ$；
书末数据相容性已单独保留，且说明小角近似限制。
